\begin{center}
	\zihao{4}\heiti
	%
	\LARGE{ 中国科学院大学\\$\boldsymbol{2014}$年招收攻读硕士学位研究生入学统一考试试卷\\
	科目名称：电动力学}

\end{center}
\noindent
考生须知：\\
1.本试卷满分为150分,全部考试时间总计180分钟。\\
2.所有答案必须写在答题纸上,写在试题纸上或草稿纸上一律无效。

\vspace*{2ex}
\hrule
\vspace*{4ex}

\fancypagestyle{mypagestyle}{%
	\fancyhf{}% Clear header/footer
	\fancyfoot[OR]{第~\thepage~页，共2页}%
	\fancyfoot[EL]{第~\thepage~页，共2页}%
	\fancyfoot[OL]{科目名称：电动力学 }%
	\fancyfoot[ER]{科目名称：电动力学 }%
	\renewcommand{\headrulewidth}{0pt}% Header rule of .4pt
	\renewcommand{\footrulewidth}{0.5pt}% Header rule of .4pt
}
\fancypagestyle{plain}{%
	\fancyhf{}% Clear header/footer
	\fancyfoot[OR]{第~\thepage~页，共2页}%
	\fancyfoot[EL]{第~\thepage~页，共2页}%
	\fancyfoot[OL]{科目名称：电动力学 }%
	\fancyfoot[ER]{科目名称：电动力学 }%
}

\setcounter{page}{1}%设置页面计数器
\pagestyle{mypagestyle}

\onehalfspacing
% 排版结束，下面是题目

\vspace{0.4cm}

\noindent
一、简答题（共24分，每题8分）
\vspace{0.2cm}

\noindent
1.\hspace{0.8em}写出真空中电磁场能量守恒定律的微分表达式，并写出相应的电磁场能量密度$\omega$及能流密度$\Vec{S}$的表达式。
\vspace{0.2cm}

\noindent
2.\hspace{0.8em}写出洛伦兹规范下电磁场矢势$\Vec{A}$和标势$\varphi$满足的达朗贝尔方程及推迟势解。
\vspace{0.2cm}

\noindent
3.\hspace{0.8em}若将四维时空矢量记为$x_{\mu} = (\Vec{x}, ict)$，写出对应的真空中的电磁场张量$F_{\mu \nu}$\vspace{0.5cm}

\noindent
二、\hspace{0.2em}单项选择题（共36分，每题6分）
\vspace{0.2cm}

\noindent
1.\hspace{0.8em}已知磁感应强度$\Vec{B} = B_0 \Vec{e_z}$，以下哪个表达式是其对应的磁场的矢势$\Vec{A}$
\vspace{0.2cm}

A.\hspace{0.8em} $\Vec{A} = (-B_0y , 0, 0)$; \hspace{5em} B.\hspace{0.8em} $\Vec{A} = (B_0y , B_0 x, 0)$; 
\vspace{0.2cm}

C.\hspace{0.8em} $\Vec{A} = (0 , -B_0 x, 0)$; \hspace{5em} D.\hspace{0.8em} $\Vec{A} = (2 B_0y , 2 B_0 x, 0)$。
\vspace{0.2cm}

\noindent
2.\hspace{0.8em}若某区域内任意一点$\Vec{r}$处的静磁场可以用磁标势描述，则其前提是
\vspace{0.2cm}

A.\hspace{0.8em}该区域内各处极化电流密度为零
\vspace{0.2cm}

B.\hspace{0.8em}$\Vec{r}$处的自由电流密度为零
\vspace{0.2cm}

C.\hspace{0.8em}电流密度满足连续性方程
\vspace{0.2cm}

D.\hspace{0.8em}$\Vec{H}$沿该区域内任何闭合回路的积分均为零
\vspace{0.3cm}

\noindent
3.\hspace{0.8em}假设海水的磁导率$\mu$可以取为真空磁导率$\mu_0 = 4 \pi \times 10^{-7} \mathrm {H} \cdot \mathrm{m^{-1} } $，电导率为$\sigma = 1 (\Omega \cdot \mathrm{m )^{-1} }$，那么频率为0.4 MHz的电磁波在海水中的穿透深度为
\vspace{0.2cm}

A.\hspace{0.8em}20m\hspace{5em}B.\hspace{0.8em}0.02m\hspace{5em}C.\hspace{0.8em}0.8m\hspace{5em}D.\hspace{0.8em}130m 
\vspace{0.3cm}

\noindent
4.\hspace{0.8em}一个矩形谐振腔的边长分别是1cm, 2cm, 3cm, 则波长是$\displaystyle{ \frac{4}{\sqrt{5}} }$cm到$\displaystyle{ \frac{8}{\sqrt{13}} }$CM之间的电磁波的模式有多少个
\vspace{0.2cm}

A.\hspace{0.8em}3个\hspace{5em}B.\hspace{0.8em}4个\hspace{5em}C.\hspace{0.8em}5个\hspace{5em}D.\hspace{0.8em}6个
\vspace{0.3cm}

\noindent
5.\hspace{0.8em}波长为$\lambda$的光入射到半径为$l$的介质微球上($l \ll \lambda $)，使微球产生辐射，下面关于其辐射总功率的说法正确的是
\vspace{0.4cm}

A.\hspace{0.8em}随着$l$的增大而增大，与$\lambda$的平方成反比\vspace{0.2cm}

B.\hspace{0.8em}随着$l$的增大而减小，与$\lambda$的平方成正比\vspace{0.2cm}

C.\hspace{0.8em}随着$l$的增大而增大，与$\lambda$的四次方成反比\vspace{0.2cm}

D.\hspace{0.8em}随着$l$的增大而减小，与$\lambda$的四次方成正比

\newpage

\noindent
6.\hspace{0.8em}一把直尺相对于坐标系$\Sigma$静止，直尺与$x$轴的夹角为$\theta$。一观察者以速度$v$沿$x$轴运动，他看到直尺与$x$轴的夹角$\theta^{'}$变为
\vspace{0.3cm}

A.\hspace{0.8em} $\sin{\theta^{'}} = \displaystyle{ \frac{\sin \theta}{ \sqrt{1 - v^2 / c^2  }}}$\hspace{8em}
B.\hspace{0.8em} $\tan{\theta^{'}} = \displaystyle{ \frac{\tan \theta}{ \sqrt{1 - v^2 / c^2  }}}$\vspace{0.3cm}

C.\hspace{0.8em}$\tan{\theta^{'}} = \displaystyle{  \sqrt{1 - v^2 / c^2  } \tan \theta } $\hspace{6.6em}
D.\hspace{0.8em} $\theta^{'} = \theta $\vspace{0.4cm}

\noindent
三、\hspace{0.2em}（共30分）在介电常数为$\varepsilon_1$，半径为$R$的均匀介质球的中心放置一个电偶极子$\Vec{p_f} = p_f \Vec{e_z}$，球外为真空。
\vspace{0.2cm}

\noindent
1. 写出介质球球面处电势满足的边界条件;

\noindent
2. 利用分离变量法求出球内外空间各点的电势;

\noindent
3. 求出介质球面上的极化面电荷密度。\vspace{3cm}

\noindent
四、\hspace{0.2em}（共30分）有三层平行介质，介电常数分别为$\varepsilon_1$，$\varepsilon_2$，$\varepsilon_3$，其中$\varepsilon_3 = \varepsilon_1$，且磁导率相同，即$\mu_1 = \mu_2 = \mu_3$。介质1、3皆可视为无限厚，中间介质2的厚度为d，现有一平面电磁波从介质1垂直于介质1、2的界面入射。
\vspace{0.2cm}

\noindent
1.列出两个界面处电磁波满足的边界条件;

\noindent
2.求出在介质1中观察到的电场强度的反射系数;

\noindent
3.根据以上结果，说明什么条件下反射系数等于0（即完全透射）。
\vspace{3cm}

\noindent
五、\hspace{0.5em}（共30分）如图所示有一小型空心带电介质圆盘，内半径为$R_1$，外半径为$R_2$，盘的厚度可以忽略不计，盘上带有均匀电荷，面电荷密度为$\sigma$，圆盘以恒定角速度$\omega$绕过其中心的垂直轴转动。设该中垂线为$x$轴，轴上$A$点与圆盘中心$O$之间的距离为$x$ ($x \gg R_2$)。
\vspace{0.2cm}

\noindent
1. 求带电空心圆盘的磁偶极距;

\noindent
2. 求A点处的磁感应强度;

\noindent
3. 如果圆盘转动角速度不是定值而是$\omega = \omega_0 \sin \alpha t $，$\alpha$是一个小常数($R_2 \ll x \ll c/\alpha$， c为光速)。在A点处平行于圆盘放一半径为b的足够小导体线圈，圆心在A点处。求出线圈中感生电场的强度。

\vspace{-0.8cm}
\begin{figure}[h!]
    \flushright
    \includegraphics[height=5cm]{test/2014/pic_1.png}
\end{figure}